\documentclass{article}
\usepackage[czech]{babel}
\usepackage[utf8]{inputenc}
\usepackage{amsmath}
\usepackage{enumitem}

\renewcommand{\div}{\operatorname{div}}
\newcommand{\R}{\mathbb R}
\newcommand{\vc}[1]{\boldsymbol{#1}}
\def\d{{\rm{d}}}

\begin{document}

\title{Test č. 2 -- F}
\date{12. 12. 2024}
\author{}
\maketitle


\noindent
Určete, zda posloupnosti funkcí

\begin{enumerate}[label=(\alph*)]
\item
\begin{align*}
        f_n(x) = \frac{\cos (nx)}{\sqrt{n}}, \qquad x\in[0,2\pi],
\end{align*}
%konverguje ke své bodové limitě v normě $L_2$ a $L_\infty$.

\item
\begin{align*}
        g_n(x) = |x|^{\frac{1}{n}}, \qquad x\in[-1,1],
\end{align*}
\end{enumerate}

konvergují ke své bodové limitě v normě $L_2$ a $L_\infty$.

\vspace{2cm}

Postup:
\begin{enumerate}
\item Načtrněte graf několika prvních funkcí posloupnosti.
\item Určete bodovou limitu $f$, $g$.
\item Zapište symbolicky, co znamená, že $f_n$ konverguje k $f$ ve zvolené normě, resp $g_n$ konverguje k $g$.
\item Ověřte výpočtem.
\end{enumerate}


\vfill
Pomůcka k integraci:
\[\cos^2(\theta) = \frac{1 + \cos(2\theta)}{2}\]

% \vfill

\end{document}
